%% =====================================================================
%%  B5-T-07-01 -- what B5 contributes to "Listening as an Instrument"
%%  -------------------------------------------------------------------
%%  LAYER A: APPEND-ONLY. Written once, then left alone. Material found
%%  only here sits in the `unique` slot and is protected by rule R10.
%% =====================================================================
%% @kind: source
%% @id: B5-T-07-01
%% @book: B5
%% @topic: T-07-01
%% @locator: ch. 2, pp. 36--69
%% @status: distilled
%% @unique: yes
%% @integrated: yes
%% @updated: 2026-09-19

\dossier{B5}{T-07-01}{\bookshort{B5}}{ch. 2, pp. 36--69}

\begin{distilled}
  \item The source contributes a named delivery channel for listening-driven work: the late-night FM DJ voice --- \emph{deep, soft, slow, and reassuring} --- credited in the opening case with de-escalating a confrontational call before any technique was applied.
  \item Voice is presented as a switch the negotiator sets deliberately: DJ voice for de-escalation, normal voice for rapport, assertive voice reserved and mostly discouraged.
  \item The chapter's overall doctrine is that the counterpart's words are only part of the channel; how you sound while collecting them decides whether collection is permitted.
  \item It treats listening itself as the offensive instrument: the counterpart who feels heard keeps supplying the information that moves the case.
\end{distilled}

\begin{unique}
  \item The voice-down doctrine as a distinct, practised skill --- not a personality trait but a settable setting, named and drilled.
  \item The pairing rule: collection techniques (mirrors, labels, silence) fail in an assertive voice; the delivery channel is a precondition, not a decoration.
\end{unique}
